% Options for packages loaded elsewhere
\PassOptionsToPackage{unicode}{hyperref}
\PassOptionsToPackage{hyphens}{url}
\documentclass[
]{article}
\usepackage{xcolor}
\usepackage[margin=1in]{geometry}
\usepackage{amsmath,amssymb}
\setcounter{secnumdepth}{5}
\usepackage{iftex}
\ifPDFTeX
  \usepackage[T1]{fontenc}
  \usepackage[utf8]{inputenc}
  \usepackage{textcomp} % provide euro and other symbols
\else % if luatex or xetex
  \usepackage{unicode-math} % this also loads fontspec
  \defaultfontfeatures{Scale=MatchLowercase}
  \defaultfontfeatures[\rmfamily]{Ligatures=TeX,Scale=1}
\fi
\usepackage{lmodern}
\ifPDFTeX\else
  % xetex/luatex font selection
\fi
% Use upquote if available, for straight quotes in verbatim environments
\IfFileExists{upquote.sty}{\usepackage{upquote}}{}
\IfFileExists{microtype.sty}{% use microtype if available
  \usepackage[]{microtype}
  \UseMicrotypeSet[protrusion]{basicmath} % disable protrusion for tt fonts
}{}
\makeatletter
\@ifundefined{KOMAClassName}{% if non-KOMA class
  \IfFileExists{parskip.sty}{%
    \usepackage{parskip}
  }{% else
    \setlength{\parindent}{0pt}
    \setlength{\parskip}{6pt plus 2pt minus 1pt}}
}{% if KOMA class
  \KOMAoptions{parskip=half}}
\makeatother
\usepackage{color}
\usepackage{fancyvrb}
\newcommand{\VerbBar}{|}
\newcommand{\VERB}{\Verb[commandchars=\\\{\}]}
\DefineVerbatimEnvironment{Highlighting}{Verbatim}{commandchars=\\\{\}}
% Add ',fontsize=\small' for more characters per line
\usepackage{framed}
\definecolor{shadecolor}{RGB}{248,248,248}
\newenvironment{Shaded}{\begin{snugshade}}{\end{snugshade}}
\newcommand{\AlertTok}[1]{\textcolor[rgb]{0.94,0.16,0.16}{#1}}
\newcommand{\AnnotationTok}[1]{\textcolor[rgb]{0.56,0.35,0.01}{\textbf{\textit{#1}}}}
\newcommand{\AttributeTok}[1]{\textcolor[rgb]{0.13,0.29,0.53}{#1}}
\newcommand{\BaseNTok}[1]{\textcolor[rgb]{0.00,0.00,0.81}{#1}}
\newcommand{\BuiltInTok}[1]{#1}
\newcommand{\CharTok}[1]{\textcolor[rgb]{0.31,0.60,0.02}{#1}}
\newcommand{\CommentTok}[1]{\textcolor[rgb]{0.56,0.35,0.01}{\textit{#1}}}
\newcommand{\CommentVarTok}[1]{\textcolor[rgb]{0.56,0.35,0.01}{\textbf{\textit{#1}}}}
\newcommand{\ConstantTok}[1]{\textcolor[rgb]{0.56,0.35,0.01}{#1}}
\newcommand{\ControlFlowTok}[1]{\textcolor[rgb]{0.13,0.29,0.53}{\textbf{#1}}}
\newcommand{\DataTypeTok}[1]{\textcolor[rgb]{0.13,0.29,0.53}{#1}}
\newcommand{\DecValTok}[1]{\textcolor[rgb]{0.00,0.00,0.81}{#1}}
\newcommand{\DocumentationTok}[1]{\textcolor[rgb]{0.56,0.35,0.01}{\textbf{\textit{#1}}}}
\newcommand{\ErrorTok}[1]{\textcolor[rgb]{0.64,0.00,0.00}{\textbf{#1}}}
\newcommand{\ExtensionTok}[1]{#1}
\newcommand{\FloatTok}[1]{\textcolor[rgb]{0.00,0.00,0.81}{#1}}
\newcommand{\FunctionTok}[1]{\textcolor[rgb]{0.13,0.29,0.53}{\textbf{#1}}}
\newcommand{\ImportTok}[1]{#1}
\newcommand{\InformationTok}[1]{\textcolor[rgb]{0.56,0.35,0.01}{\textbf{\textit{#1}}}}
\newcommand{\KeywordTok}[1]{\textcolor[rgb]{0.13,0.29,0.53}{\textbf{#1}}}
\newcommand{\NormalTok}[1]{#1}
\newcommand{\OperatorTok}[1]{\textcolor[rgb]{0.81,0.36,0.00}{\textbf{#1}}}
\newcommand{\OtherTok}[1]{\textcolor[rgb]{0.56,0.35,0.01}{#1}}
\newcommand{\PreprocessorTok}[1]{\textcolor[rgb]{0.56,0.35,0.01}{\textit{#1}}}
\newcommand{\RegionMarkerTok}[1]{#1}
\newcommand{\SpecialCharTok}[1]{\textcolor[rgb]{0.81,0.36,0.00}{\textbf{#1}}}
\newcommand{\SpecialStringTok}[1]{\textcolor[rgb]{0.31,0.60,0.02}{#1}}
\newcommand{\StringTok}[1]{\textcolor[rgb]{0.31,0.60,0.02}{#1}}
\newcommand{\VariableTok}[1]{\textcolor[rgb]{0.00,0.00,0.00}{#1}}
\newcommand{\VerbatimStringTok}[1]{\textcolor[rgb]{0.31,0.60,0.02}{#1}}
\newcommand{\WarningTok}[1]{\textcolor[rgb]{0.56,0.35,0.01}{\textbf{\textit{#1}}}}
\usepackage{graphicx}
\makeatletter
\newsavebox\pandoc@box
\newcommand*\pandocbounded[1]{% scales image to fit in text height/width
  \sbox\pandoc@box{#1}%
  \Gscale@div\@tempa{\textheight}{\dimexpr\ht\pandoc@box+\dp\pandoc@box\relax}%
  \Gscale@div\@tempb{\linewidth}{\wd\pandoc@box}%
  \ifdim\@tempb\p@<\@tempa\p@\let\@tempa\@tempb\fi% select the smaller of both
  \ifdim\@tempa\p@<\p@\scalebox{\@tempa}{\usebox\pandoc@box}%
  \else\usebox{\pandoc@box}%
  \fi%
}
% Set default figure placement to htbp
\def\fps@figure{htbp}
\makeatother
\setlength{\emergencystretch}{3em} % prevent overfull lines
\providecommand{\tightlist}{%
  \setlength{\itemsep}{0pt}\setlength{\parskip}{0pt}}
\usepackage{bookmark}
\IfFileExists{xurl.sty}{\usepackage{xurl}}{} % add URL line breaks if available
\urlstyle{same}
\hypersetup{
  pdftitle={Modelação Ecológica - Tutorial Aula 1 - Final Task},
  pdfauthor={Tiago A. Marques},
  hidelinks,
  pdfcreator={LaTeX via pandoc}}

\title{Modelação Ecológica - Tutorial Aula 1 - Final Task}
\author{Tiago A. Marques}
\date{October 25, 2018}

\begin{document}
\maketitle

\tableofcontents

\section{Introduction}\label{introduction}

This document describes a possible implementation of the final task
presented to the students of ``Modelação Ecológica'' under the tutorial
``A hands-on tutorial on R and R Studio''.

\section{The task}\label{the-task}

Here we will implement an exercise were we pretend we are sampling an
animal population, using some (very basic) simulations to understand the
process better. Create plots that represent all the steps of your task,
with proper legends, labels, colors, etc, and add all your comments to
the dynamic report. 1. Simulate the positions of 10000 animals in a
study area, with length 10km and width 1km. Assume that any animal has
an equal chance to be at any location in the study area (this
corresponds to a uniform density surface). 2. Generate a transect at a
random location along the study area. 3. Assume that you can potentially
detect at most animals up to 500 meters from your transect. Count all
the animals that you would detect if detection was perfect across your
transect. 4. Consider that animals far from your transect are harder to
detect - yes, you are doing distance sampling! Define a function that
represents a distance sampling half-normal detection function (If
needed, look back at your slides from the 1st class). Assume that
sigma=200m. 5. Simulate the detection process and get a sample of those
animals detected. 6. Create a plot that allows you to estimate (at this
stage just a visual guess in needed) the detection probability. 7.
Repeat the sampling process 500 times, and store the number of animals
detected in each one of your simulated surveys. 8. plot the distribution
of the number of animals that you would detect each new survey. Take
your own conclusions about all that you did.

\section{Implementation}\label{implementation}

\subsection{Simulate animal positions}\label{simulate-animal-positions}

First we define some state variables

\begin{Shaded}
\begin{Highlighting}[]
\CommentTok{\#number of animals}
\NormalTok{N}\OtherTok{=}\DecValTok{10000}
\CommentTok{\#units in meters}
\CommentTok{\#length}
\NormalTok{L}\OtherTok{=}\DecValTok{10000}
\CommentTok{\#width}
\NormalTok{W}\OtherTok{=}\DecValTok{1000}
\end{Highlighting}
\end{Shaded}

Next, we simulate the animals, and show them in a plot. Since the
distribution was uniform in space, we use the function \texttt{runif}

\begin{Shaded}
\begin{Highlighting}[]
\CommentTok{\#set the seed to get always the same results}
\FunctionTok{set.seed}\NormalTok{(}\DecValTok{123}\NormalTok{)}
\NormalTok{xs}\OtherTok{=}\FunctionTok{runif}\NormalTok{(N,}\DecValTok{0}\NormalTok{,W)}
\NormalTok{ys}\OtherTok{=}\FunctionTok{runif}\NormalTok{(N,}\DecValTok{0}\NormalTok{,L)}
\FunctionTok{par}\NormalTok{(}\AttributeTok{mfrow=}\FunctionTok{c}\NormalTok{(}\DecValTok{1}\NormalTok{,}\DecValTok{1}\NormalTok{))}
\FunctionTok{plot}\NormalTok{(xs,ys,}\AttributeTok{pch=}\StringTok{"."}\NormalTok{,}\AttributeTok{main=}\StringTok{"The simulated animal locations"}\NormalTok{)}
\end{Highlighting}
\end{Shaded}

\pandocbounded{\includegraphics[keepaspectratio]{FinalTask_files/FinalTask_files/figure-latex/unnamed-chunk-2-1.pdf}}

\subsection{Generate a transect}\label{generate-a-transect}

Next, we need to generate the transect. The transect goes along the
study area. After the fact I realized this is not sensible, since the
area is only 1km wide and we will assume that we could see up to 500
meters on each side. Therefore, on the fly, I am going to change the
width of the area to be 20 km, so that the exercise is more sensible.

The new area and animals are shown below:

\begin{Shaded}
\begin{Highlighting}[]
\CommentTok{\#number of animals}
\NormalTok{N}\OtherTok{=}\DecValTok{10000}
\CommentTok{\#units in meters}
\CommentTok{\#length}
\NormalTok{L}\OtherTok{=}\DecValTok{10000}
\CommentTok{\#width}
\NormalTok{W}\OtherTok{=}\DecValTok{20000}
\CommentTok{\#set the seed to get always the same results}
\FunctionTok{set.seed}\NormalTok{(}\DecValTok{123}\NormalTok{)}
\NormalTok{xs}\OtherTok{=}\FunctionTok{runif}\NormalTok{(N,}\DecValTok{0}\NormalTok{,W)}
\NormalTok{ys}\OtherTok{=}\FunctionTok{runif}\NormalTok{(N,}\DecValTok{0}\NormalTok{,L)}
\FunctionTok{par}\NormalTok{(}\AttributeTok{mfrow=}\FunctionTok{c}\NormalTok{(}\DecValTok{1}\NormalTok{,}\DecValTok{1}\NormalTok{))}
\FunctionTok{plot}\NormalTok{(xs,ys,}\AttributeTok{pch=}\StringTok{"."}\NormalTok{,}\AttributeTok{main=}\StringTok{"The simulated animal locations"}\NormalTok{)}
\end{Highlighting}
\end{Shaded}

\pandocbounded{\includegraphics[keepaspectratio]{FinalTask_files/FinalTask_files/figure-latex/unnamed-chunk-3-1.pdf}}

Then, I generate a transect along the study area, and to avoid edge
effects, I assume the transect will not overlap the edges (by defining a
buffer where the transect midle point can not be placed).

\begin{Shaded}
\begin{Highlighting}[]
\CommentTok{\#define a buffer}
\NormalTok{buffer}\OtherTok{=}\DecValTok{500}
\CommentTok{\#get transect location}
\NormalTok{trans}\OtherTok{=}\FunctionTok{runif}\NormalTok{(}\DecValTok{1}\NormalTok{,}\DecValTok{0}\SpecialCharTok{+}\NormalTok{buffer,W}\SpecialCharTok{{-}}\NormalTok{buffer)}
\FunctionTok{par}\NormalTok{(}\AttributeTok{mfrow=}\FunctionTok{c}\NormalTok{(}\DecValTok{1}\NormalTok{,}\DecValTok{1}\NormalTok{))}
\FunctionTok{plot}\NormalTok{(xs,ys,}\AttributeTok{pch=}\StringTok{"."}\NormalTok{,}\AttributeTok{main=}\StringTok{"Simulated animal locations"}\NormalTok{)}
\FunctionTok{abline}\NormalTok{(}\AttributeTok{v=}\NormalTok{trans,}\AttributeTok{lwd=}\DecValTok{2}\NormalTok{,}\AttributeTok{col=}\DecValTok{4}\NormalTok{)}
\end{Highlighting}
\end{Shaded}

\pandocbounded{\includegraphics[keepaspectratio]{FinalTask_files/FinalTask_files/figure-latex/unnamed-chunk-4-1.pdf}}

\subsection{Animals that can be
detected}\label{animals-that-can-be-detected}

We simply need to see how many animals are within 500 meters of our
transect

\begin{Shaded}
\begin{Highlighting}[]
\CommentTok{\#define what we call in distance sampling the truncation distance}
\CommentTok{\#no animals can be detected beyhond that distance}
\NormalTok{w}\OtherTok{=}\DecValTok{500}
\CommentTok{\#calculate the perpendicular distance to the transect, for each animal}
\NormalTok{dist}\OtherTok{=}\FunctionTok{abs}\NormalTok{(trans}\SpecialCharTok{{-}}\NormalTok{xs)}
\CommentTok{\#how many are at risk}
\NormalTok{natrisk}\OtherTok{=}\FunctionTok{sum}\NormalTok{(dist}\SpecialCharTok{\textless{}}\NormalTok{w)}
\end{Highlighting}
\end{Shaded}

There are 477 animals within 500

We can plot them with a different color in our study area

\begin{Shaded}
\begin{Highlighting}[]
\FunctionTok{plot}\NormalTok{(xs,ys,}\AttributeTok{col=}\FunctionTok{ifelse}\NormalTok{(dist}\SpecialCharTok{\textless{}}\NormalTok{w,}\DecValTok{2}\NormalTok{,}\DecValTok{1}\NormalTok{))}
\FunctionTok{abline}\NormalTok{(}\AttributeTok{v=}\NormalTok{trans,}\AttributeTok{lwd=}\DecValTok{2}\NormalTok{,}\AttributeTok{col=}\DecValTok{4}\NormalTok{)}
\end{Highlighting}
\end{Shaded}

\pandocbounded{\includegraphics[keepaspectratio]{FinalTask_files/FinalTask_files/figure-latex/unnamed-chunk-6-1.pdf}}

\subsection{Defining a detection
function}\label{defining-a-detection-function}

Using the material from the slides, we can define the detection
function. Below we plot the detection function setting
\texttt{sigma=200} by default

\begin{Shaded}
\begin{Highlighting}[]
\NormalTok{HN}\OtherTok{=}\ControlFlowTok{function}\NormalTok{(x,}\AttributeTok{sigma=}\DecValTok{200}\NormalTok{)\{}
\NormalTok{  px}\OtherTok{=}\FunctionTok{exp}\NormalTok{(}\SpecialCharTok{{-}}\NormalTok{(x}\SpecialCharTok{\^{}}\DecValTok{2}\NormalTok{)}\SpecialCharTok{/}\NormalTok{(}\DecValTok{2}\SpecialCharTok{*}\NormalTok{sigma}\SpecialCharTok{\^{}}\DecValTok{2}\NormalTok{))}
  \FunctionTok{return}\NormalTok{(px)}
\NormalTok{\}}
\FunctionTok{plot}\NormalTok{(}\DecValTok{0}\SpecialCharTok{:}\DecValTok{500}\NormalTok{,}\FunctionTok{HN}\NormalTok{(}\DecValTok{0}\SpecialCharTok{:}\DecValTok{500}\NormalTok{),}\AttributeTok{xlab=}\StringTok{"x=distance (m)"}\NormalTok{,}\AttributeTok{ylab=}\StringTok{"Detection probability g(x)"}\NormalTok{)}
\end{Highlighting}
\end{Shaded}

\pandocbounded{\includegraphics[keepaspectratio]{FinalTask_files/FinalTask_files/figure-latex/unnamed-chunk-7-1.pdf}}

\subsection{Simulate the detection
process}\label{simulate-the-detection-process}

Now we can simulate the detection process. To do that we first evaluate
the detection probability of each animal, and then compare that
probability with a random draw from a Uniform in (0,1). If the value is
smaller than the probability of detection, the animal is detected, if
not, it is missed

\begin{Shaded}
\begin{Highlighting}[]
\CommentTok{\#P detectar cada um dos meus individuos}
\NormalTok{pdet}\OtherTok{=}\FunctionTok{HN}\NormalTok{(dist)}
\CommentTok{\#an index for animals potentially detected}
\NormalTok{inTRANSECT}\OtherTok{=}\NormalTok{dist}\SpecialCharTok{\textless{}}\DecValTok{500}
\CommentTok{\#get coordinates of just those}
\NormalTok{xs2}\OtherTok{=}\NormalTok{xs[inTRANSECT]}
\NormalTok{ys2}\OtherTok{=}\NormalTok{ys[inTRANSECT]}
\CommentTok{\#their distances}
\NormalTok{dist}\OtherTok{=}\NormalTok{dist[inTRANSECT]}
\CommentTok{\#and their pdets}
\NormalTok{pdet}\OtherTok{=}\NormalTok{pdet[inTRANSECT]}
\CommentTok{\#how }
\NormalTok{n1}\OtherTok{=}\FunctionTok{sum}\NormalTok{(inTRANSECT)}
\CommentTok{\#get the uniform}
\NormalTok{punif}\OtherTok{=}\FunctionTok{runif}\NormalTok{(n1,}\DecValTok{0}\NormalTok{,}\DecValTok{1}\NormalTok{)}
\CommentTok{\#index for those detected detected}
\NormalTok{inDET}\OtherTok{=}\NormalTok{pdet}\SpecialCharTok{\textgreater{}}\NormalTok{punif}
\CommentTok{\#how many where they?}
\NormalTok{n2}\OtherTok{=}\FunctionTok{sum}\NormalTok{(inDET)}
\end{Highlighting}
\end{Shaded}

Now, we can plot the detected animals on top of the previous plots,
zooming in the sampled transect

\begin{Shaded}
\begin{Highlighting}[]
\FunctionTok{plot}\NormalTok{(xs,ys,}\AttributeTok{xlim=}\FunctionTok{c}\NormalTok{(trans}\SpecialCharTok{{-}}\NormalTok{(w}\SpecialCharTok{+}\DecValTok{50}\NormalTok{),trans}\SpecialCharTok{+}\NormalTok{(w}\SpecialCharTok{+}\DecValTok{50}\NormalTok{)))}
\FunctionTok{points}\NormalTok{(xs2,ys2,}\AttributeTok{col=}\StringTok{"red"}\NormalTok{,}\AttributeTok{bg=}\StringTok{"red"}\NormalTok{,}\AttributeTok{pch=}\DecValTok{21}\NormalTok{)}
\FunctionTok{abline}\NormalTok{(}\AttributeTok{v=}\NormalTok{trans,}\AttributeTok{lwd=}\DecValTok{2}\NormalTok{,}\AttributeTok{col=}\DecValTok{4}\NormalTok{)}
\FunctionTok{points}\NormalTok{(xs2[inDET],ys2[inDET],}\AttributeTok{col=}\StringTok{"green"}\NormalTok{,}\AttributeTok{bg=}\StringTok{"blue"}\NormalTok{,}\AttributeTok{pch=}\DecValTok{21}\NormalTok{)}
\end{Highlighting}
\end{Shaded}

\pandocbounded{\includegraphics[keepaspectratio]{FinalTask_files/FinalTask_files/figure-latex/unnamed-chunk-9-1.pdf}}

We had 477 animals potentially detected, and we detected 237.

\subsection{The detection probability}\label{the-detection-probability}

We simply need an histogram of the distances of animals at risk of being
detected, and we can add on top the detected distances

\begin{Shaded}
\begin{Highlighting}[]
\FunctionTok{hist}\NormalTok{(dist,}\AttributeTok{main=}\StringTok{"Detected distances"}\NormalTok{,}\AttributeTok{xlab=}\StringTok{"Distance (m)"}\NormalTok{,}\AttributeTok{ylab=}\StringTok{"g(x)"}\NormalTok{)}
\FunctionTok{hist}\NormalTok{(dist[inDET],}\AttributeTok{add=}\ConstantTok{TRUE}\NormalTok{,}\AttributeTok{col=}\DecValTok{3}\NormalTok{)}
\end{Highlighting}
\end{Shaded}

\pandocbounded{\includegraphics[keepaspectratio]{FinalTask_files/FinalTask_files/figure-latex/unnamed-chunk-10-1.pdf}}

Roughtly looking at this, we would estimate a probability of detection
of about 50 \%.

We actually detected 49.7 \% of the animals in the covered transect. Not
bad for a ``guessestimate'' above!

Note that we can actually calculate the true probability of detection,
by integrating the detection function (the area under the curve) and
dividing that by \texttt{w} (the area of the histogram if we detected
all distances).

\begin{Shaded}
\begin{Highlighting}[]
\NormalTok{trueP}\OtherTok{=}\FunctionTok{integrate}\NormalTok{(}\AttributeTok{f =}\NormalTok{ HN,}\AttributeTok{lower =} \DecValTok{0}\NormalTok{,}\AttributeTok{upper=}\NormalTok{w)}\SpecialCharTok{$}\NormalTok{value}\SpecialCharTok{/}\NormalTok{w}
\end{Highlighting}
\end{Shaded}

The true probability of detecting an animal within 500m of the transect
is 0.495.

\subsection{Repeat the process many
times}\label{repeat-the-process-many-times}

We could now simply repeat all this many times. Note that we could
repeat the process of generating the animals' positions, or the
positions and the transect, or condition on the observed animals and
transect, and then repeat just the detection component. Here we do the
latter, but leave the former as an exercise for you to do at home.

\begin{Shaded}
\begin{Highlighting}[]
\CommentTok{\#how many times do we repeat the process}
\NormalTok{n3}\OtherTok{=}\DecValTok{200}
\NormalTok{numberdetected}\OtherTok{=}\FunctionTok{numeric}\NormalTok{(n3)}
\ControlFlowTok{for}\NormalTok{(j }\ControlFlowTok{in} \DecValTok{1}\SpecialCharTok{:}\NormalTok{n3)\{}
\CommentTok{\#get the uniform}
\NormalTok{punif}\OtherTok{=}\FunctionTok{runif}\NormalTok{(n1,}\DecValTok{0}\NormalTok{,}\DecValTok{1}\NormalTok{)}
\CommentTok{\#index for those detected detected}
\NormalTok{inDET}\OtherTok{=}\NormalTok{pdet}\SpecialCharTok{\textgreater{}}\NormalTok{punif}
\CommentTok{\#how many where they?}
\NormalTok{numberdetected[j]}\OtherTok{=}\FunctionTok{sum}\NormalTok{(inDET)}
\NormalTok{\}}
\end{Highlighting}
\end{Shaded}

\subsection{Plot the outcome}\label{plot-the-outcome}

Finally, we could plot the number of animals detected in each of our 200
simulations

\begin{Shaded}
\begin{Highlighting}[]
\FunctionTok{hist}\NormalTok{(numberdetected)}
\end{Highlighting}
\end{Shaded}

\pandocbounded{\includegraphics[keepaspectratio]{FinalTask_files/FinalTask_files/figure-latex/unnamed-chunk-13-1.pdf}}

\section{Conclusions}\label{conclusions}

We were able to simulate an animal population and a simple sampling
strategy, to formulate a model for detectability (based on distance
sampling) and to simulate the detection process.

We investigated the randomness involved in the number of animals
detected conditional on the animals available and a transect position,
and noted there was considerable variability across simulations.

\end{document}
